\chapter*{Mở đầu}
\addcontentsline{toc}{chapter}{Mở đầu}
\markboth{MỞ ĐẦU}{}

\section*{Lý do chọn đề tài}

Trong những năm gần đây, mô hình ngôn ngữ lớn (Large Language Model --- LLM) được tích hợp ngày càng sâu vào các hệ thống truy vấn tri thức doanh nghiệp, trợ lý tra cứu cơ sở dữ liệu (Text-to-SQL), tác tử đa công cụ (Multi-Tool Agent) và kiến trúc sinh văn bản tăng cường truy hồi (Retrieval-Augmented Generation --- RAG). Lợi ích về khả năng hiểu ngôn ngữ tự nhiên đi kèm các rủi ro an toàn đặc thù tại thời điểm chạy (runtime): tấn công chèn chỉ dẫn trực tiếp (direct prompt injection), tiêm nhiễm gián tiếp qua tài liệu hoặc ô dữ liệu trong cơ sở dữ liệu (indirect prompt/SQL injection), lạm dụng công cụ (tool misuse) và rò rỉ dữ liệu nhạy cảm. Việc chỉ dựa vào khả năng ``tự căn chỉnh'' (alignment) của mô hình hoặc chỉ kiểm thử khép kín trên tập dữ liệu tự sinh là không đủ; một cơ chế phòng thủ thực tiễn đòi hỏi kiến trúc nhiều lớp theo chiều sâu (defense-in-depth), kiểm soát truy cập trước truy hồi, đồng thời phải được kiểm chứng độc lập bằng các bộ dữ liệu công khai, công cụ red-team chuẩn và đo đạc chi phí độ trễ định lượng.

\section*{Bối cảnh và vấn đề}

Hệ thống RAG và Agent doanh nghiệp kết hợp kho tài liệu không cấu trúc với cơ sở dữ liệu có cấu trúc (SQL) và các công cụ thực thi (gửi email, tạo phiếu yêu cầu), tạo ra bốn bề mặt tấn công liên thông: (1)~đầu vào người dùng cố ghi đè chỉ dẫn hệ thống hoặc mã hoá payload qua Base64/ROT13/Leetspeak; (2)~tài liệu trong kho tri thức hoặc ô văn bản trong bảng SQL bị đầu độc chỉ thị ẩn; (3)~câu lệnh SQL do mô hình sinh ra cố vượt quyền đọc bảng/cột nhạy cảm; và (4)~đầu ra mô hình vô tình phát tán bí mật, PII hoặc liên kết thất thoát dữ liệu. Bên cạnh đó, hạn chế phổ biến của nhiều nghiên cứu quy mô phòng thí nghiệm là tình trạng \emph{đánh giá khép kín} (tự sinh mẫu tấn công, tự viết luật regex, tự chấm điểm), thiếu đo đạc phân vị độ trễ (\texttt{p50/\allowbreak p95/\allowbreak p99}) và thiếu đối sánh định lượng với các giải pháp chuẩn quốc tế như Meta Llama Guard~3 hay NVIDIA NeMo Guardrails. Đề tài Đồ án tốt nghiệp này giải quyết trực diện các khoảng trống đó.

\section*{Mục tiêu}

\begin{enumerate}
    \item Nghiên cứu toàn diện bề mặt tấn công của hệ thống RAG, Text-to-SQL và Multi-Tool Agent theo khung OWASP Top~10 for LLM Applications; khảo sát và phân tích cơ chế của các giải pháp phòng thủ chuẩn (\texttt{Llama Guard~3}, \texttt{NeMo Guardrails}) cùng các khung kiểm thử đối kháng độc lập (\texttt{NVIDIA garak}, \texttt{Microsoft PyRIT}).
    \item Thiết kế và hiện thực khung bảo mật đa lớp \textbf{Shield AI}: kết hợp tầng tiền giải mã (\emph{Payload De-obfuscation Pre-decoder} cho Base64/ROT13/Leetspeak/Unicode), chuỗi Guard tất định (Input Guard, Provenance Guard, RAG Context Guard, Output/DLP Guard), bộ phán đoán ngữ nghĩa cục bộ (\emph{Semantic Judge} \texttt{gemma4:e2b-\allowbreak it-\allowbreak q4\_\allowbreak K\_\allowbreak M}) và cơ chế kiểm soát truy cập theo vai trò (\emph{Pre-retrieval RBAC}) trên nền truy hồi lai vectorized.
    \item Mở rộng kiến trúc bảo vệ sang dữ liệu có cấu trúc và tác tử đa công cụ: xây dựng \textbf{SQL Guard} dùng chính trình biên dịch \texttt{EXPLAIN} và \texttt{Authorizer AST callback} của SQLite để kiểm soát quyền truy cập tới từng bảng/cột, \textbf{Indirect Result Guard} quét ô dữ liệu trả về từ CSDL, và \textbf{Multi-Tool Company Agent} có hàng chờ phê duyệt cấp quản lý (\emph{Leader Approval Queue}).
    \item Thực hiện chương trình kiểm thử đối kháng và đánh giá đối sánh độc lập theo bốn trục: (i)~đánh giá trên benchmark công khai ngoài phân phối \texttt{deepset/\allowbreak prompt-\allowbreak injections} ($662$ mẫu đa ngôn ngữ); (ii)~bắn trực tiếp công cụ \texttt{NVIDIA garak v0.\allowbreak 17.\allowbreak 1} ($2.048$ probe qua HTTP REST) và bộ biến đổi/tấn công thích nghi đa lượt \texttt{PyRIT} vào endpoint hệ thống; (iii)~đo đạc chi tiết phân vị độ trễ \texttt{p50/\allowbreak p95/\allowbreak p99} của từng tầng xử lý khi bật/tắt Semantic Judge; và (iv)~lập bảng đối sánh định lượng công bằng với \texttt{Meta Llama Guard~3 1B} và \texttt{NVIDIA NeMo Guardrails}.
\end{enumerate}

\section*{Đối tượng và phạm vi}

\textbf{Đối tượng} nghiên cứu là kiến trúc Guardrails đa lớp kết hợp kiểm soát truy cập RBAC cho hệ thống RAG, Text-to-SQL và Multi-Tool Agent, cùng phương pháp luận kiểm thử đối kháng độc lập và đánh giá đối sánh định lượng. \textbf{Phạm vi} thực nghiệm sử dụng bộ dữ liệu doanh nghiệp tổng hợp đa phòng ban (gồm kho tri thức 48 tài liệu của hệ thống RAG và cơ sở dữ liệu 12 bảng / 29 tài liệu của doanh nghiệp nội thất Nhà Mộc Việt), kết hợp với bộ benchmark công khai \texttt{deepset/\allowbreak prompt-\allowbreak injections} từ HuggingFace và bộ probe gốc của \texttt{garak v0.\allowbreak 17.\allowbreak 1}; toàn bộ suy luận LLM và Semantic Judge chạy cục bộ qua \texttt{Ollama} trên GPU cá nhân (RTX~4050 6\,GB VRAM) để đảm bảo tính tái lập 100\% và không phát sinh chi phí API thương mại.

\section*{Câu hỏi nghiên cứu}

Chương trình thực nghiệm của đồ án (Chương~4) được thiết kế để trả lời năm câu hỏi nghiên cứu định lượng:

\begin{itemize}
    \item \textbf{CH1 (Độ chính xác và khái quát hoá ngoài phân phối).} Bộ luật tất định đạt hiệu quả như thế nào trên tập kiểm thử nội bộ ($425$ case) so với khi kiểm chứng độc lập trên benchmark công khai ngoài phân phối \texttt{deepset/\allowbreak prompt-\allowbreak injections} ($662$ mẫu đa ngôn ngữ), và việc bổ sung tầng tiền giải mã (\emph{Pre-decoder}) cùng luật cấu trúc đa ngôn ngữ cải thiện khả năng khái quát hoá ra sao mà vẫn giữ $\text{FPR} \approx 0\%$?
    \item \textbf{CH2 (Đối sánh với giải pháp chuẩn quốc tế).} Khi chạy trên cùng tập kiểm thử và cùng phần cứng cục bộ, kiến trúc lai của \textbf{Shield AI} (Rule + Pre-decoder + Semantic Judge \texttt{gemma4:e2b}) đạt chỉ số $\text{TPR}$, $\text{FPR}$, $\text{F1}$ và độ trễ như thế nào so với \textbf{Meta Llama Guard~3 1B} và \textbf{NVIDIA NeMo Guardrails} (\texttt{self\_\allowbreak check\_\allowbreak input})?
    \item \textbf{CH3 (Phân tích độ trễ từng tầng p50/p95/p99).} Chi phí độ trễ (\texttt{mean/\allowbreak p50/\allowbreak p95/\allowbreak p99}) của từng lớp Guard tất định và của toàn chuỗi khi bật/tắt Semantic Judge là bao nhiêu, và cơ chế lọc phân tầng (\emph{tiered short-circuit}) giúp giảm tỷ lệ phải gọi LLM Judge như thế nào?
    \item \textbf{CH4 (Chống chịu trước công cụ Red-Team độc lập \texttt{garak} và \texttt{PyRIT} thích nghi).} Khi bắn trực tiếp \texttt{NVIDIA garak v0.\allowbreak 17.\allowbreak 1} ($1.024$ probe trên mỗi endpoint) và chiến dịch \texttt{PyRIT} (35 biến thể mã hoá Converter + Attacker LLM tự động viết lại prompt qua 4 lượt hội thoại) vào API sống, tỷ lệ ngăn chặn ở từng lớp đạt bao nhiêu và cơ chế \emph{Pre-retrieval RBAC + Defense-in-Depth} bảo vệ dữ liệu nhạy cảm như thế nào khi kẻ tấn công thích nghi vượt qua lớp lọc đầu vào?
    \item \textbf{CH5 (Bảo vệ dữ liệu có cấu trúc Text-to-SQL và Multi-Tool Agent).} Cơ chế \textbf{SQL Guard} dựa trên trình biên dịch SQLite \texttt{EXPLAIN} + Authorizer AST và \textbf{Indirect Result Guard} đạt độ chính xác ($\text{TPR}/\text{FPR}/\text{F1}$) cùng độ trễ \texttt{p50/\allowbreak p95/\allowbreak p99} như thế nào trước các đòn tấn công P2SQL trực tiếp ($69$ mẫu) và tiêm nhiễm gián tiếp qua ô dữ liệu DB ($70$ mẫu dev/holdout)?
\end{itemize}

\section*{Phương pháp}

Nhóm áp dụng phương pháp nghiên cứu thực nghiệm có kiểm soát: (1)~khảo sát chuẩn OWASP LLM Top~10, OWASP LLMSVS, NIST AI~600-1 và các công trình học thuật gốc; (2)~thiết kế kiến trúc phòng thủ đa lớp fail-closed tách biệt rõ lớp kiểm soát quyền truy cập tất định (RBAC/SQL Authorizer) với lớp kiểm duyệt nội dung lai (Regex + Pre-decoder + Semantic Judge); (3)~kiểm chứng chéo bằng công cụ độc lập (\texttt{garak v0.\allowbreak 17.\allowbreak 1}, \texttt{PyRIT}, \texttt{deepset/\allowbreak prompt-\allowbreak injections}) và đối sánh baseline (\texttt{Llama Guard~3}, \texttt{NeMo Guardrails}); (4)~khóa mã băm SHA-256 cho toàn bộ tập dữ liệu và artifact kết quả (\texttt{summary.\allowbreak json}, \texttt{result.\allowbreak jsonl}) để mọi con số trong báo cáo đều tái lập được từ mã nguồn.

\section*{Đóng góp chính của Đồ án tốt nghiệp}

\begin{itemize}
    \item \textbf{Kiến trúc Guardrails lai đa lớp tích hợp Tiền giải mã (Pre-decoder) và Pre-retrieval RBAC:} Kết hợp giải mã tự động Base64/ROT13/Leetspeak/Unicode ở tầng tiền xử lý, bộ luật cấu trúc đa ngôn ngữ, kiểm soát truy cập RBAC trước truy hồi, RAG Context Guard, Output/DLP Guard và Semantic Judge cục bộ (\texttt{gemma4:e2b-\allowbreak it-\allowbreak q4\_\allowbreak K\_\allowbreak M}).
    \item \textbf{Mở rộng bảo vệ sang Text-to-SQL (SQLite AST Authorizer) và Multi-Tool Company Agent:} Dùng chính trình biên dịch \texttt{EXPLAIN} và callback \texttt{set\_\allowbreak authorizer} của SQLite để phân tích chính xác bảng/cột được đọc (bắt trọn CTE, subquery, \texttt{UNION}, \texttt{SELECT *}), kết hợp \textbf{Indirect Result Guard} quét từng ô dữ liệu trả về và hàng chờ phê duyệt hành động nhạy cảm (\emph{Leader Approval Queue}) trên bộ dữ liệu doanh nghiệp Nhà Mộc Việt (12 bảng SQL, 469 unit/integration tests PASS).
    \item \textbf{Đo đạc phân vị độ trễ p50/p95/p99 toàn diện và tối ưu hoá Semantic Judge:} Phát hiện và khắc phục triệt để lỗi sinh chuỗi suy nghĩ ngầm (\texttt{think=False}) giúp giảm độ trễ gọi Judge từ $\sim 6.440\,\text{ms}$ xuống $\text{p50} = 99{,}79\,\text{ms}$; chứng minh tổng chi phí của toàn bộ 6 lớp Guard tất định chỉ tốn $\mathbf{\text{p50} = 0{,}42\,\text{ms}, \text{p95} = 1{,}21\,\text{ms}, \text{p99} = 1{,}39\,\text{ms}}$ (trên $1.275$ lượt đo).
    \item \textbf{Kiểm chứng độc lập trên Benchmark công khai và đối sánh định lượng với Llama Guard~3 \& NeMo Guardrails:} Chứng minh định lượng hiện tượng khớp quá mức (overfitting) của bộ regex cũ trên \texttt{deepset/\allowbreak prompt-\allowbreak injections} ($662$ mẫu: TPR rớt từ $97{,}0\%$ xuống $4{,}2\%$), và sau khi nâng cấp kiến trúc P5 đạt \textbf{$\text{TPR} = 68{,}1\%$ (Rule-Only, $\text{FPR} = 0{,}0\%$)} và \textbf{$\text{TPR} = 87{,}5\%$, $\text{F1} = 92{,}9\%$ (Rule + Judge, $\text{FPR} = 0{,}5\%$)} --- vượt xa \textbf{Meta Llama Guard~3 1B} ($\text{F1} = 36{,}9\%$) và \textbf{NVIDIA NeMo Guardrails} ($\text{F1} = 45{,}2\%$).
    \item \textbf{Kiểm thử đối kháng hộp đen trực tiếp bằng \texttt{NVIDIA garak v0.\allowbreak 17.\allowbreak 1} và \texttt{PyRIT} Adaptive Multi-Turn:} Chặn thành công \textbf{$\mathbf{950/1.024}$ probe ($\mathbf{92{,}8\%}$)} của \texttt{garak} trên endpoint \texttt{/\allowbreak v1/\allowbreak rag/\allowbreak query} (trong đó \texttt{HijackHateHumans} đạt $100\%$, \texttt{LatentJailbreak} đạt $100\%$, \texttt{InjectBase64} đạt $97{,}7\%$), chặn \textbf{$\mathbf{35/35}$ ($\mathbf{100\%}$)} biến thể mã hoá của \texttt{PyRIT Converter Matrix}, và giữ tỷ lệ rò rỉ dữ liệu mục tiêu ở mức \textbf{$\mathbf{0/5}$ ($\mathbf{0\%}$)} trước kẻ tấn công thích nghi đa lượt nhờ kiến trúc phòng thủ chiều sâu.
\end{itemize}
